\section{Experimental evaluation}
\label{sec:submission-results}
\subsection{Cohort and protocol}
The primary experiment uses a fixed cohort of 39 graphs. Each graph is
evaluated with $(h,k)=(1,1)$, $(2,1)$ and $(3,2)$, giving 117 labeling instances.
It comprises nine random labeled trees with 20, 40 and 60 vertices and
three graph-generation seeds per size; three grids
$P_4\square P_5$, $P_5\square P_8$ and $P_6\square P_{10}$;
18 Erd\H{o}s--R\'enyi graphs with the same sizes and seeds and edge
probabilities $p\in\{0.15,0.30\}$; and nine Barab\'asi--Albert graphs
with the same sizes and seeds and attachment parameter $m=2$.
Disconnected samples are retained. The graph cohort is not reselected
according to solver outcomes.

The SAT backend is CaDiCaL with order encoding and a fresh solver for each
span trial. The ILP baseline is Gurobi with the assignment formulation,
one thread and a greedy MIP start. Symmetry breaking is disabled in this
comparison. Each instance has three repetitions per backend, producing
\ProgressObservations{} observations. Backends run sequentially in
separate processes, with their order rotated across instances and
repetitions. Each run has a 30-second external limit; elapsed wall time
includes startup, model construction, solving and cleanup. Parent-side
witness validation is outside the recorded worker interval.
The shared comparison runner's external deadline governs both backends;
the individual solver APIs' time-limit mechanisms are not assumed equivalent.

The completed \texttt{isolated-progress-v2} run retains intermediate
validated incumbents and lower-bound updates. SAT transmits completed
span decisions; Gurobi transmits available incumbent and bound information.
A timeout without a closed interval is reported as \FEAS, not \OPT.
One Gurobi observation returned \FEAS{} early; its native stopping reason
was not recorded, so it remains a feasible observation with an unknown
native termination reason. No \OPT{} observation exceeds the recorded budget.

The source manifest, explicit graphs, configurations, bounds and saved
labelings are retained with the raw results. An audit checks all
\ProgressObservations{} witnesses and the expected cohort and repetitions.
Label validation independently checks shortest-path constraints, but does
not verify SAT UNSAT traces or MILP dual certificates. Historical pilot,
v1 and supplementary $L(2,1)$ symmetry results remain in the companion
research report and archive; their timings are not pooled into this experiment.

\subsection{Optimality coverage and bound quality}
\begin{table}[htbp]
\centering\small
\caption{Optimality coverage within 30 seconds.}
\label{tab:submission-coverage}
\input{\ProgressReportDir/coverage.tex}
\tablenote{Each backend has 351 observations on 117 instances. All-three OPT
counts instances proved optimal in every repetition. Improved UB counts
FEASIBLE observations with a saved feasible span smaller than the shared
initial incumbent. It does not establish optimality.}
\end{table}

Table~\ref{tab:submission-coverage} shows more stable optimality proofs for
Gurobi: \ProgressILPStable/117 instances, compared with
\ProgressSATStable/117 for CaDiCaL. Both backends are optimal in all three
repetitions on \ProgressJointStable{} instances. Repetitions are observations
of the same graphs, not additional independent graph samples. Runtime
comparisons restricted to this common solved subset would need to be read
alongside the coverage difference; they cannot characterize unsolved instances.

\begin{table}[htbp]
\centering\small
\caption{Paired lower- and upper-bound comparisons.}
\label{tab:submission-bounds}
\input{\ProgressReportDir/bounds.tex}
\tablenote{Pairs match the same graph, parameter pair and repetition. A
higher lower bound and a smaller validated upper bound are preferable.
The open subset comprises the 21 instances unresolved after pooling v2
solver evidence, before either combinatorial postprocessing step. Its
63 pairs are retained as a fixed analysis subset; they are not the final
\ConstructionOpen{} open instances in Table~\ref{tab:submission-evidence}.}
\end{table}

On this initially open subset, SAT has a smaller feasible upper bound in
\ProgressOpenSATUB{} pairs, ILP in \ProgressOpenILPUB{}, and
\ProgressOpenEqualUB{} tie. ILP has a stronger recorded lower bound in
\ProgressOpenILPLB{} pairs and \ProgressOpenEqualLB{} tie; SAT has none
with a stronger lower bound on this subset. These outcomes distinguish
solution quality from the ability to exclude smaller spans. They support
complementary behavior on the measured instances, not equivalent
effectiveness or universal dominance by either method.

\subsection{Consolidated optimality evidence}
\begin{table}[htbp]
\centering\small
\caption{Consolidated instance-level evidence.}
\label{tab:submission-evidence}
\input{\ConstructionReportDir/summary.tex}
\tablenote{Original-v2 columns: Theory denotes a witness attaining the
applicable initial theoretical bound; Other denotes additional closed
intervals. Theory + Other + Open = Cases. The final columns include both
certification steps: two new closures from cliques with stored witnesses
and three from color-order constructions. Clique-cert. counts matching
square-clique bounds and labelings; it overlaps Closed and is not added
again. Closed + Open = Cases. Counts are distinct graph/parameter
instances, not per-run successes; timed solver coverage remains unchanged.}
\end{table}

Applicable formulas and bounds provide a separate correctness reference.
In particular, the equal-parameter tree formula and the $L(2,1)$ grid
formula in the appendix cover 12 instances in this cohort. References
are used only within their stated graph and parameter domains. Agreement
on tested instances is validation evidence, rather than a proof of a new
formula over an entire family.

Pooling v2 solver intervals originally closed \ProgressClosed{} instances.
The square-clique analysis increased this to \EvidenceClosedAfter{} by
matching a lower bound to stored feasible witnesses in two previously
unresolved cases. The additional construction in
Section~\ref{sec:color-order-results} supplies new explicit witnesses and
closes \ConstructionNewlyClosed{} more. Table~\ref{tab:submission-evidence}
is the final instance-level summary. Solver-dependent proofs remain
distinguished from directly checked combinatorial certificates.

A separate screen of \TreeGraphCount{} $L(3,2)$ trees with 100--1600 vertices
found matching reported optimal spans for both backends; \TreeTightCount{}
attain the degree bound. Its different size range, single run per graph
and backend, and 60-second limit prevent pooling its timings with the
primary experiment. It supplies additional structured-graph validation,
not evidence that every tree is easy or that the degree bound is always tight.

\subsection{Validity and reproducibility limits}
The cohort includes several graph models but only three graph-generation
seeds at each randomized configuration and 20--60 vertices in the primary
comparison. Its observations do not represent all graphs or all solver
implementations. Assignment ILP is one formulation; the results are not a
comparison with every possible ILP method. SAT/ILP agreement supplements,
but does not replace, mathematical arguments or independent certificates.

The experiment does not include a controlled direct-versus-order encoding
ablation, incremental SAT or a GA baseline. Clause counts measure model
representation and do not prove runtime improvements. The post-hoc
construction is analyzed separately from timed solver performance.
Historical hardware metadata are incomplete, so no reference-machine
performance claim is inferred from the available software versions alone.

Raw observations, witnesses and manifests remain under \path{results/runs/}.
The v2 analysis, combinatorial certificates and consolidated summaries are
stored in versioned directories under \path{results/analysis/}. Source
fingerprints link derived tables to their inputs. The report-generation
targets revalidate saved evidence and compile the manuscript without
rerunning the original solver experiment.
